\documentclass[../main.tex]{subfiles}

\begin{document}

\section*{Kết quả đạt được}

Khóa luận đã xây dựng được hệ thống quản lý đồ án tốt nghiệp cho Khoa Công nghệ thông tin, chạy trên máy chủ thật qua HTTPS, bao phủ quy trình từ phân công giảng viên hướng dẫn tới lưu kho đồ án. Đối chiếu với bốn mục tiêu trong đề cương, các kết quả cụ thể như sau.

\textbf{Mục tiêu 1 -- Khảo sát, phân tích quy trình.} Quy trình quản lý đồ án tại Khoa được khảo sát qua phiếu hỏi 18 sinh viên, các biểu mẫu thực tế (mẫu đề cương H01, bảng định hướng đề tài, quy tắc đặt tên tệp) và trao đổi với giảng viên hướng dẫn. Kết quả phân tích gồm sáu cặp quy trình AS-IS và TO-BE mô hình hóa bằng BPMN, 49 yêu cầu chức năng thuộc tám nhóm, 19 yêu cầu phi chức năng có ngưỡng đo được, sơ đồ use case với năm use case chính cùng nhóm chức năng chung và đặc tả chi tiết 29 use case con (chương~2). Quá trình khảo sát cũng chỉ ra những điểm đề cương cần điều chỉnh cho khớp với thực tế, như việc giảng viên giao đề tài thay cho sinh viên tự đăng ký và quy trình duyệt đề cương hai cấp (giảng viên hướng dẫn duyệt, sau đó Lãnh đạo Bộ môn ký duyệt).

\textbf{Mục tiêu 2 -- Thiết kế và xây dựng hệ thống theo kiến trúc SOA.} Hệ thống gồm ba thành phần độc lập: client Web viết bằng React, Web API viết bằng FastAPI và dịch vụ AI tách riêng, giao tiếp qua REST và chỉ có một cửa vào công khai là Apache với HTTPS. Mô hình miền gồm 22 lớp thực thể; CSDL MySQL chia thành hai schema do hai dịch vụ sở hữu riêng, các quy tắc nghiệp vụ quan trọng được bảo đảm bằng ràng buộc của CSDL (chương~3). Năm use case chính đã được thi công đầy đủ cho ba vai trò sinh viên, giảng viên và Bộ môn: quản lý đề tài gắn với định hướng, đề cương nộp nhiều phiên bản và được duyệt hai cấp (giảng viên hướng dẫn, sau đó Bộ môn phê duyệt bằng bản PDF đề cương đã ký), báo cáo tiến độ theo mốc, hội đồng và lịch bảo vệ có chấm và chốt điểm, kho đồ án có import CSV và tra cứu. Hệ thống cài đặt và nâng cấp bằng một lệnh nhờ Docker Compose.

\textbf{Mục tiêu 3 -- Tích hợp AI kiểm tra trùng lặp.} Dịch vụ AI dùng mô hình \texttt{vietnamese-bi-encoder} để so sánh tên đề tài theo ngữ nghĩa ngay khi giảng viên tạo hoặc giao đề tài, và kết hợp vân tay văn bản theo cụm âm tiết (bắt sao chép nguyên văn, kể cả khi sửa vài từ, bỏ qua cụm câu khuôn mẫu xuất hiện ở nhiều đề tài) với embedding (bắt diễn đạt lại) để phát hiện các đoạn trùng trong đề cương, báo cáo định dạng DOCX và PDF, chỉ rõ đồ án nguồn và từng đoạn trùng. Tài liệu bị cảnh báo khi trùng nhiều với một đồ án, hoặc khi tổng tỷ lệ trùng với mọi đồ án vượt ngưỡng tổng, để bắt cả bài ghép từ nhiều nguồn. Mô hình được chọn bằng một quy tắc viết trước khi đo; ngưỡng được hiệu chỉnh trên một tập dữ liệu và kiểm chứng trên tập giữ riêng. Trên tập kiểm chứng, độ chính xác ở mức tài liệu là 7/8 cặp (0,875), ở mức tên đề tài là 22/24 truy vấn (0,917); mọi cặp chép nguyên văn đều bị cảnh báo và mọi cặp không liên quan đều dưới ngưỡng. Toàn bộ việc kiểm tra chạy trên máy chủ của Khoa, nội dung đồ án không gửi ra dịch vụ bên ngoài.

\textbf{Mục tiêu 4 -- Triển khai thử nghiệm và đánh giá.} Hệ thống đã được triển khai trên máy chủ thật và nạp dữ liệu thử nghiệm vượt các chỉ tiêu đặt ra: 580 tài khoản sinh viên, 70 tài khoản giảng viên, 300 đề tài trong kho, 80 truy vấn có nhãn và 50 tệp tài liệu gồm 25 cặp. Kho trên máy chủ cũng đã được nạp 106 trong 128 đồ án thật đã thu thập của khoa. Kiểm thử tự động gồm 71 ca cho Web API, 15 ca cho dịch vụ AI cùng các kịch bản E2E và đo đạc qua HTTPS. Đối chiếu với 19 yêu cầu phi chức năng, 17 yêu cầu đạt hoàn toàn, hai yêu cầu đạt một phần (phân quyền chưa chạy hết kịch bản E2E của UC2--UC5, giao diện mới kiểm thử tự động trên Chrome). Các chỉ tiêu thời gian đạt với khoảng dư lớn: kiểm tra tên đề tài qua HTTPS có P95 dưới 0,12 giây so với chỉ tiêu 5 giây, kiểm tra một đề cương có P95 dưới 5 giây so với chỉ tiêu 30 giây (chương~4).

Như vậy, các mục tiêu đề ra trong đề cương về cơ bản đã được hoàn thành. Hệ thống đáp ứng các vấn đề mà sinh viên nêu ra trong khảo sát: theo dõi được trạng thái đề tài, đề cương trên một nơi chung; phát hiện sớm đề tài, nội dung trùng với các đồ án đã có ngay từ khâu giao đề tài và duyệt đề cương; lưu lại đồ án các khóa để tra cứu.

\section*{Hạn chế}

Bên cạnh các kết quả trên, khóa luận còn một số hạn chế:
\begin{itemize}
    \item \textit{Về khảo sát:} phiếu khảo sát mới thu được ý kiến của 18 sinh viên; quy trình phía giảng viên và Bộ môn được tìm hiểu qua biểu mẫu và trao đổi trực tiếp, chưa có khảo sát định lượng với các nhóm này.
    \item \textit{Về dữ liệu đánh giá AI:} độ chính xác kiểm tra trùng lặp được đo trên dữ liệu tổng hợp do bộ sinh dữ liệu tạo ra, vì kho đồ án thật chưa được gán nhãn. Phát hiện câu diễn đạt lại còn yếu (40--56\% số câu), và khi kho có nhiều đồ án cùng chủ đề vẫn có thể cảnh báo nhầm. Ngưỡng tổng mới chỉ được kiểm tra trên các bài ghép dựng thêm, chưa có bài ghép thật. Báo cáo rất dài cần thời gian kiểm tra tăng theo độ dài (khoảng một phút cho 350\,000 ký tự). Kho thật trên máy chủ còn 22 đồ án chưa nạp được do chỉ có tệp nguồn LaTeX.
    \item \textit{Về phê duyệt đề cương:} hệ thống chỉ yêu cầu bản phê duyệt là tệp PDF, không xác thực chữ ký trên đề cương (dù là chữ ký tay đã quét hay chữ ký điện tử); việc bảo đảm bản tải lên là bản đã ký thuộc về người phê duyệt.
    \item \textit{Về chức năng và bảo mật:} người dùng chưa tự đổi được mật khẩu; hệ thống chưa có sao lưu CSDL tự động; tài liệu OpenAPI vẫn truy cập được khi chưa đăng nhập. Hệ thống chưa có đăng ký đề tài trực tuyến, trao đổi trực tiếp giữa sinh viên và giảng viên hay nhắc hạn tự động, là những tính năng sinh viên có đề xuất trong khảo sát.
    \item \textit{Về triển khai và kiểm thử:} hệ thống chạy trên một máy chủ, không có dự phòng; giao diện chỉ thiết kế cho máy tính, một số bảng như \textit{Tiến độ toàn khoa} còn hẹp cột tên đề tài. Kịch bản E2E của UC2--UC5 và kịch bản chụp màn hình chưa được cập nhật theo luồng duyệt đề cương hai cấp; giao diện chưa được kiểm thử tự động trên Edge; trên máy chủ thật mới chạy các phép đo không ghi dữ liệu. Hệ thống chưa được dùng cho một đợt đồ án thật với người dùng thật.
\end{itemize}

\section*{Hướng phát triển}

Từ các hạn chế trên, hệ thống có thể được phát triển tiếp theo các hướng sau:
\begin{itemize}
    \item \textit{Đưa vào sử dụng thử:} vận hành hệ thống cho một đợt đồ án của Khoa, thu thập ý kiến của sinh viên, giảng viên và Bộ môn để điều chỉnh quy trình và giao diện.
    \item \textit{Nâng cao chất lượng kiểm tra trùng lặp:} gán nhãn một phần kho đồ án thật để hiệu chỉnh lại ngưỡng; thử các mô hình hoặc bước xếp hạng lại (cross-encoder) cho câu diễn đạt lại; hỗ trợ trích xuất văn bản từ mã nguồn LaTeX để nạp nốt các đồ án còn thiếu; chia nhỏ việc kiểm tra báo cáo dài để rút ngắn thời gian chờ.
    \item \textit{Xác thực chữ ký khi phê duyệt:} nếu Khoa thống nhất dùng chữ ký số cho đề cương, bổ sung bước xác minh chữ ký và chuỗi chứng thư tới các tổ chức cung cấp dịch vụ chứng thực chữ ký số được công nhận, như một tùy chọn bên cạnh bản ký tay đã quét.
    \item \textit{Bổ sung chức năng:} cho người dùng tự đổi và khôi phục mật khẩu, đăng nhập bằng tài khoản của trường; đăng ký đề tài trực tuyến theo danh sách đề tài giảng viên đưa ra; nhắc hạn nộp qua email; trao đổi giữa sinh viên và giảng viên ngay trên đề tài; thống kê, báo cáo theo khóa, theo giảng viên.
    \item \textit{Hoàn thiện vận hành:} sao lưu CSDL và tệp tải lên tự động; giới hạn truy cập tài liệu OpenAPI; cập nhật các kịch bản kiểm thử E2E và tích hợp chúng vào quy trình kiểm thử tự động trước mỗi lần nâng cấp; bổ sung giao diện cho thiết bị di động.
\end{itemize}

\end{document}
